\chapter{Conclusions}
\label{chap:conclusions}

% ─────────────────────────────────────────────────────────────────────────────
\section{Summary}
\label{sec:summary}

This thesis presented Pseudo++, an open-source toolchain for the Romanian
baccalaureate pseudocode dialect.  The motivating gap was clear: despite
pseudocode being a central part of the Romanian high-school CS curriculum,
no complete execution environment existed for it.  Students had to
mentally trace programs by hand, with no automated way to verify whether
a solution was correct.

The toolchain addresses this gap through eight components built around a
single C core library: a Tree-sitter grammar that formally captures the
dialect's syntax, a source formatter that normalises Unicode notation to
an ASCII form the parser can process, a tree-walking interpreter, a
step-through debugger with live variable and condition inspection, a
two-pass transpiler targeting C, C++, and Pascal, a loop-equivalence
transformation module, a browser-based web editor that requires no
installation and runs entirely offline via WebAssembly, and a
command-line interface for scripting and batch use.

A deliberate architectural decision runs through all components: the
same C library is compiled both to a native binary and to a WebAssembly
module, so the CLI and the web editor exercise identical logic.  There is
no server involved in execution --- the entire toolchain runs in the
browser once the page is loaded.

% ─────────────────────────────────────────────────────────────────────────────
\section{Limitations}
\label{sec:limitations}

Several limitations are worth acknowledging.

\paragraph{Scalar variables only.}
The language processed by Pseudo++ covers the subset of the baccalaureate
pseudocode that uses scalar variables (integers, reals, strings).  Arrays
and records are not part of this dialect as taught at the secondary level,
so they are outside the scope of the toolchain.  Programs that require
indexed data structures cannot be expressed or executed.

\paragraph{Forward-only debugging.}
The debugger supports stepping forward through execution one statement at
a time but does not support stepping backwards.  Inspecting earlier
program states requires restarting execution from the beginning.

\paragraph{Heuristic re-indentation.}
The formatter's structural re-indentation pass is keyword-driven: it
increases or decreases indentation based on the presence of specific
tokens at line boundaries.  Pathological source formatting --- deeply
mixed Unicode and ASCII notation, or unusual whitespace structure ---
can cause the pass to produce incorrect indentation, leading to parse
errors for programs that would otherwise be syntactically valid.

\paragraph{Limited static analysis.}
The \texttt{check} module reports a basic set of semantic warnings.
It does not perform data-flow analysis, dead-code detection, or
type-mismatch warnings across branches, all of which would be useful
for an educational linter.

% ─────────────────────────────────────────────────────────────────────────────
\section{Future Work}
\label{sec:future}

\paragraph{A dialect-agnostic architecture via a common IR.}
The current toolchain is tightly coupled to the baccalaureate dialect:
the interpreter and transpiler operate directly on the Tree-sitter CST,
so supporting a new dialect would require duplicating significant backend
logic.  A more scalable direction is to introduce a dialect-agnostic
intermediate representation (IR) as a common target.  Each dialect would
need only a grammar and a front-end pass that lowers its CST to the IR;
the interpreter, transpiler, debugger, and static analyser would then
operate entirely at the IR level and be reusable across all dialects.
This architecture would make it straightforward to add support for other
Romanian pseudocode variants --- such as the dialect used in the
Babeș-Bolyai University admission examination, which extends the
baccalaureate notation with arrays, subprograms, and additional
built-in operations --- without touching any backend code.

\paragraph{Richer static analysis.}
Expanding the \texttt{check} module with data-flow analysis would enable
diagnostics such as use-before-assignment, unreachable code after
unconditional branches, and type inconsistencies across conditional
paths.  These are common sources of logic errors in student code and
are straightforwardly detectable on a CST with a simple fixed-point
analysis.

\paragraph{Step-back debugging.}
Implementing a snapshot-based mechanism that captures interpreter state
before each step would enable backward navigation through execution
history.  Given the explicit-stack architecture of the interpreter,
snapshots reduce to copying the stack, environment, and a small set of
auxiliary fields --- a feasible extension that would bring the debugging
experience closer to omniscient debuggers studied in the literature
\cite{lewis2003omniscient,reversible-debugging}.

\paragraph{Integration with educational platforms.}
Pseudo++'s WebAssembly module could be integrated into existing Romanian
educational platforms such as pbinfo.ro, either as an inline widget for
running pseudocode examples or as a fully supported submission language,
allowing students to solve problems directly in pseudocode without first
translating to C++ or Pascal.
